\documentclass{SibFU-docs}

% доп. пакеты
\usepackage{graphicx}
\usepackage{float}
\usepackage{hyperref}
\usepackage{caption}
\usepackage{subcaption}
\usepackage{amsmath}
% предотвартим ошибку с \Bbbk перед загрузкой пакета amssymb
\let\Bbbk\relax
\usepackage{amssymb}
\usepackage{booktabs}
\usepackage{tabularx}
\usepackage{longtable}
\usepackage{enumitem}
\usepackage{multirow}

% гиперссылки
\hypersetup{
    colorlinks=true,
    linkcolor=black,
    citecolor=blue,
    urlcolor=blue
}


% библиография
\addbibresource{report.bib}

% титульный лист
\begin{document}
\setlist[itemize]{left=1.3cm, itemsep=1pt, topsep=2pt} % для выравнивания списков
\setlist[enumerate]{left=1.3cm, itemsep=1pt, topsep=2pt}
\renewcommand{\contentsname}{}
\mycovertitle
{Гуманитарный институт}
{Кафедра прикладной информатики в искусстве и интерактивном медиа}
{Технологии создания и обработки табличных данных}
{ст. преподаватель И.~Р.~Нигматуллин}
{Захаров И.~М., Мельников С.~В., Фиряго О.~А., Тугаринова А.~И.,\\
Тетерина П.~А.}
 
\begin{center}
\bfseries РЕФЕРАТ
\end{center}
\noindent

Реферат по теме: «Технологии создания и обработки табличных данных» содержит 28 страниц и 37 использованных источников.

Ключевые слова: ТАБЛИЧНЫЕ ДАННЫЕ, ЭЛЕКТРОННЫЕ ТАБЛИЦЫ, СТРУКТУРА ДАННЫХ, ФОРМАТЫ ХРАНЕНИЯ (CSV, XLSX, JSON, Parquet), АНАЛИЗ ДАННЫХ, ВИЗУАЛИЗАЦИЯ, ФОРМУЛЫ И ФУНКЦИИ, АВТОМАТИЗАЦИЯ РАСЧЁТОВ, ПОДГОТОВКА И ОЧИСТКА ДАННЫХ, ИССЛЕДОВАТЕЛЬСКИЙ ИНСТРУМЕНТАРИЙ

В работе проводится комплексное исследование современных технологий и методологий работы с табличными данными — от их базовой структуризации до сложного анализа и визуального представления.

Рассмотрены фундаментальные принципы организации «чистых» данных в электронных таблицах, обеспечивающие эффективность последующей обработки. Детально проанализированы основные форматы хранения и обмена табличной информации — от универсальных текстовых (CSV, JSON) и офисных (XLSX) до специализированных высокопроизводительных (Parquet, Avro), — с выделением их специфики, преимуществ и областей применения. Особое внимание уделено инструментарию анализа (сводные таблицы, динамические массивы) и ключевым принципам эффективной визуализации для преобразования данных в инсайты.

Завершающая часть работы посвящена технологиям автоматизации расчётов с помощью формул и функций, процессу подготовки данных к анализу и роли табличных процессоров как универсального инструмента в учебной и научно-исследовательской деятельности.

Основные выводы:
\begin{enumerate}
\item Корректная структуризация данных на основе принципов «чистых данных» (одна строка — наблюдение, один столбец — переменная) является критическим фундаментом, предопределяющим возможность их последующего автоматизированного анализа, фильтрации и визуализации с помощью встроенных инструментов табличных процессоров;
\item Выбор формата хранения табличных данных должен носить осознанный характер, определяемый конкретной задачей: CSV и JSON — для простого обмена и веб-разработки, XLSX — для комплексных отчётов с форматированием и формулами, Parquet — для высокопроизводительной аналитики больших данных, что позволяет оптимизировать производительность, совместимость и затраты на хранение;
\item Современные табличные процессоры представляют собой мощные аналитические платформы, где эффективность анализа достигается за счёт синергии инструментов агрегации (сводные таблицы), динамических вычислений (формулы и функции) и визуализации, построенной на принципах информационного дизайна (минимализм, семантическое использование цвета, избегание искажений);
\item Технологии автоматизации (формулы, функции, Power Query) и методичная подготовка (очистка, преобразование) данных трансформируют табличный редактор из инструмента для простого учета в среду для воспроизводимого исследовательского цикла, что делает его незаменимым инструментом в образовании и науке для обработки данных, проверки гипотез и представления результатов.
\end{enumerate}

\vspace{1em}

Рекомендации:
\begin{itemize}[label=-]
\item При начале работы с любым массивом табличных данных в первую очередь приводить его к виду, соответствующему принципам «чистых данных», исключая объединённые ячейки, пустые строки и разнородную информацию в одном столбце, что является основой для эффективной работы;
\item Выбирать формат для долгосрочного хранения и обмена данными, учитывая требования к совместимости, функциональности и объёму: для архивов и аналитических задач отдавать предпочтение открытым и эффективным форматам, таким как XLSX (для документов) и Parquet (для больших данных);
\item Осваивать и активно использовать не только базовые, но и продвинутые инструменты анализа (сводные таблицы, функции ВПР, ФИЛЬТР, ЕСЛИ) и визуализации, соблюдая принципы чёткого и честного представления данных для повышения убедительности и ясности отчётов;
\item Внедрять в практику учебной и исследовательской работы использование инструментов автоматизированной подготовки данных (например, Power Query) и шаблонов с настроенными формулами для типовых расчетов, что экономит время, снижает количество ошибок и обеспечивает воспроизводимость результатов.
\end{itemize}

\newpage

\thispagestyle{empty}
\begin{center}
\bfseries СОДЕРЖАНИЕ
\end{center}
\vspace{-6pt}
\tableofcontents
\clearpage

\section*{}
\vspace*{-2.5em}
\begin{center}\textbf{ВВЕДЕНИЕ}\end{center}
\phantomsection
\addcontentsline{toc}{section}{Введение}

Табличные данные стали «кровеносной системой» информационного обмена в современном мире, лежа в основе отчётности, научных исследований, финансовых моделей и бизнес-аналитики. От простых списков в Excel до многомерных датасетов в облачных хранилищ — технологии создания и обработки таблиц формируют базовую цифровую грамотность и являются ключевым инструментом для преобразования сырой информации в обоснованные решения. Однако кажущаяся простота табличного интерфейса зачастую маскирует сложность грамотной работы с данными, требующей системного подхода на всех этапах — от первичной организации до финальной интерпретации.

В условиях экспоненциального роста объёмов и источников данных особенно остро встаёт проблема их качества, согласованности и пригодности для анализа. Неструктурированные, «грязные» таблицы, ошибки в формулах, неадекватный выбор формата хранения или типа визуализации приводят к искажению информации, принятию ошибочных решений и значительным временным затратам на переделку работы. Более того, даже корректно собранные данные теряют свою ценность без применения современных методов анализа и наглядного представления. В этом контексте целостное понимание технологического стека работы с табличными данными — от структурных принципов до продвинутых аналитических возможностей — становится критически важным компетенцией.

Современное состояние вопроса характеризуется богатым набором инструментов: от классических принципов структурирования («чистые данные») и множества форматов файлов (CSV, XLSX, Parquet) до мощного встроенного аналитического аппарата (функции, сводные таблицы, Power Query) в таких распространённых средах, как Microsoft Excel и Google Sheets. Однако эффективное применение этих инструментов на практике требует чёткого понимания их специфики, областей применения и взаимодополняемости.

Актуальность темы обусловлена всеобщей потребностью в data-driven подходах как в бизнесе, так и в академической среде. Грамотное владение технологиями обработки табличных данных позволяет оптимизировать временные затраты, значительно повысить точность и глубину анализа, а также качество презентации результатов. Новизна работы заключается в комплексном рассмотрении всего цикла работы с табличными данными — от фундаментальных правил их организации и выбора формата хранения до практического применения сложных методов анализа, визуализации и автоматизации в учебных и профессиональных задачах.

Цель работы: исследовать совокупность технологий и методов, составляющих современный процесс создания, обработки, анализа и визуального представления табличных данных.

Задачи:
\begin{enumerate}
\item Изучить принципы структуры и организации «чистых» табличных данных в электронных таблицах;
\item Проанализировать основные форматы хранения и обмена табличной информации, определить их особенности, преимущества, недостатки и типичные сценарии использования;
\item Рассмотреть инструментарий анализа данных (сортировка, фильтрация, сводные таблицы, функции) и принципы построения эффективных визуализаций в табличных процессорах;
\item Исследовать механизмы автоматизации расчётов с помощью формул и функций, а также процесс подготовки и очистки данных для анализа;
\item Определить роль и возможности табличных процессоров как инструмента для поддержки учебной и научно-исследовательской деятельности.
\end{enumerate}

Методы: системный анализ, сравнительный анализ форматов и инструментов, обобщение и структурирование информации из профессиональной литературы, руководств и лучших практик работы с табличными процессорами.

\newpage
% вручную увеличить номер секции и записать в оглавление с номером,
% чтобы в теле не печатался автоматический номер, но в tableofcontents он остался
\refstepcounter{section}
\addcontentsline{toc}{section}{\protect\numberline{\thesection} Структура и организация табличных данных в электронных таблицах.}
\vspace*{-2.5em}
\begin{center}\textbf{\thesection\ СТРУКТУРА И ОРГАНИЗАЦИЯ ТАБЛИЧНЫХ ДАННЫХ В ЭЛЕКТРОННЫХ ТАБЛИЦАХ}\end{center}
\vspace*{-1.5em}

\ManualSubsection{Фундаментальные элементы и архитектура табличного документа}

Электронная таблица представляет собой программную среду для структурированного представления и манипуляции данными, организованными в виде двумерной матрицы. Её архитектура строится на иерархии взаимосвязанных элементов. Базовой единицей является «ячейка» --- пересечение строки (пронумерованной) и столбца (обозначенного буквами), обладающая уникальным адресом (например, C5). Ячейки служат контейнерами для данных, которыми могут быть: текстовые надписи (метки), числовые значения (включая даты и время, хранимые как числа), логические константы (ИСТИНА/ЛОЖЬ) и, что наиболее важно, «формулы», начинающиеся со знака равенства (=) и определяющие правила вычислений. Совокупность ячеек образует «лист» (worksheet) --- основное рабочее пространство. Несколько листов объединяются в «рабочую книгу» (workbook), что позволяет логически разделять различные наборы данных или этапы работы в рамках одного файла, например, хранить исходные данные, промежуточные расчёты и итоговые отчёты на разных листах \cite{walkenbach_excel_bible_2019, gormley_data_organization_2018}.

Ключевым механизмом для построения динамических вычислений являются «ссылки». Относительные ссылки (например, A1) автоматически изменяются при копировании формулы, обеспечивая её адаптацию к новому местоположению. Абсолютные ссылки (например, \$A\$1), фиксированные знаком доллара, остаются неизменными, что необходимо для ссылок на константы (например, налоговую ставку). Смешанные ссылки (A\$1 или \$A1) фиксируют только столбец или строку. Для повышения ясности сложных формул применяются «именованные диапазоны», когда группе ячеек присваивается семантически понятное имя \cite{saffr2019}.

\ManualSubsubsection{Принципы «чистых» данных как основа аналитики}

Эффективный анализ возможен только при условии корректной первоначальной организации данных, описываемой концепцией «чистых данных» (tidy data). Её основополагающие принципы, сформулированные Хэдли Уикхемом и ставшие отраслевым стандартом, включают \cite{turingway_spreadsheets, uofiowa_spreadsheet_structure}:

\begin{enumerate}
\item «Каждая переменная формирует столбец». Столбец должен содержать данные, измеряющие один и тот же атрибут для всех наблюдений (например, только «Цена» или только «Дата продажи»). Объединение нескольких атрибутов в один столбец (например, «Фамилия Имя») затрудняет фильтрацию и поиск;
\item «Каждое наблюдение формирует строку». Строка должна представлять собой уникальную запись об одном объекте или событии (один заказ, одного сотрудника, одно измерение);
\item «Каждый тип наблюдаемых единиц формирует отдельную таблицу». Данные о разных сущностях (например, о клиентах и их заказах) должны храниться на разных листах или в разных таблицах, а связываться через уникальные идентификаторы (ID).
\end{enumerate}

Следование этим принципам автоматически запрещает использование «объединённых ячеек» в теле таблицы, так как они разрушают её регулярную структуру, делая невозможной корректную сортировку, фильтрацию и построение сводных таблиц. Также недопустимы «пустые строки и столбцы» внутри массива данных, часто используемые для визуального разделения, --- они прерывают непрерывные диапазоны и приводят к ошибкам в функциях, работающих с интервалами. «Чистая» таблица является плотной матрицей, где все ячейки в границах данных заполнены или содержат осмысленные маркеры пропусков \cite{lachermarcus_oerspreadsheets_2025}.

\ManualSubsubsection{Форматы данных и правила визуального оформления}

Помимо логической структуры, критическую роль играет правильное указание «формата данных» в ячейках. Табличный процессор интерпретирует введённую информацию в зависимости от назначенного формата: числовой (целые, дробные, проценты), денежный, финансовый, дата, время, текстовый. Например, число, введённое как текст («001»), не сможет участвовать в арифметических операциях, а дата, не распознанная системой, станет обычной строкой. Правильный формат гарантирует корректность вычислений, сортировки (даты будут отсортированы хронологически, а не лексикографически) и отображения.

«Визуальное оформление» служит цели повышения читаемости и не должно конфликтовать с «чистой» структурой. Рекомендуется: использовать жирный шрифт и заливку для строки заголовков столбцов; применять выравнивание текста по левому краю, а чисел --- по правому для визуального контраста; использовать тонкие границы для разделения данных; задействовать «условное форматирование» для автоматического выделения цветом аномалий (например, отрицательных значений). Итоговые строки (итоги, субтоталы) следует отделять сверху двойной линией. Чередование заливки фона для строк («зебра») значительно снижает утомляемость глаз при работе с большими таблицами, но должно реализовываться средствами условного форматирования, а не ручным окрашиванием ячеек \cite{gormley_data_organization_2018}.

\newpage
\refstepcounter{section}
\addcontentsline{toc}{section}{\protect\numberline{\thesection}Основные форматы хранения и обмена табличной информацией}
\begin{center}\textbf{\thesection\ ОСНОВНЫЕ ФОРМАТЫ ХРАНЕНИЯ И ОБМЕНА ТАБЛИЧНОЙ ИНФОРМАЦИЕЙ}\end{center}
\vspace*{-1.5em}
\ManualSubsection{Классификация форматов и текстовые стандарты обмена (CSV, TSV, JSON, XML)}

Многообразие форматов для хранения табличных данных обусловлено различными требованиями к производительности, функциональности, совместимости и объёму. Классифицировать их можно по нескольким ключевым признакам \cite{quadratic_common_file_formats_2025}:
\begin{itemize}[label=-]
\item «По структуре хранения»: «текстовые» (human-readable), такие как CSV, JSON, и «бинарные» (binary), такие как XLSX, Parquet;
\item «По модели данных»: «построчные» (row-oriented), оптимизированные для операций с целыми записями (CSV, Avro), и «колоночные» (column-oriented), эффективные для аналитических запросов по подмножеству полей (Parquet);
\item «По поддержке схемы»: «со схемой» (schema-on-write), где структура определяется до записи (Parquet, XLSX), и «без схемы» (schema-on-read), где структура интерпретируется при чтении (CSV, JSON).
\end{itemize}

Для простого и универсального обмена данными между системами исторически используются текстовые форматы. «CSV» (Comma-Separated Values), чья де-факто спецификация изложена в RFC 4180, представляет собой простой текст, где каждая строка --- это запись, а значения разделены запятыми. Значения, содержащие разделители или переводы строк, заключаются в кавычки. «TSV» (Tab-Separated Values) использует в качестве разделителя символ табуляции, что уменьшает вероятность конфликта, так как табуляция реже встречается в данных. Преимущества этих форматов --- предельная простота, человекочитаемость, малый размер и поддержка практически любым ПО. Главные недостатки --- отсутствие типизации данных (всё --- текст), отсутствие встроенной схемы и уязвимость к ошибкам парсинга при сложном экранировании \cite{ietf_rfc4180_csv_2005}.

Для веб-приложений и конфигурационных файлов доминирующим форматом стал «JSON» (JavaScript Object Notation, RFC 8259). Он представляет данные в виде пар «ключ-значение» и массивов. Таблицу в JSON можно представить как массив объектов, где каждый объект --- строка, а его свойства --- столбцы. Это обеспечивает отличную человекочитаемость, лёгкость парсинга в JavaScript и широкую поддержку, но ведёт к избыточности из-за повторения ключей для каждой записи и не поддерживает строгую схему по умолчанию. «XML» (eXtensible Markup Language) --- более строгий и формальный формат, позволяющий определять сложные схемы данных с помощью XSD. Однако его избыточность из-за закрывающих тегов и относительная сложность парсинга делают его менее популярным для простого обмена таблицами, хотя он является основой формата XLSX \cite{json_rfc8259_2017, ecma_office_open_xml_2016}.

\ManualSubsubsection{Форматы офисных пакетов: от проприетарных бинарных к открытым XML-стандартам}

Для хранения таблиц с полной поддержкой формул, форматирования, диаграмм, макросов и нескольких листов используются форматы офисных пакетов. Исторически это был «бинарный формат XLS Microsoft Excel», чья закрытая и сложная структура (BIFF --- Binary Interchange File Format) обеспечивала высокую производительность в самой программе, но создавала проблемы с совместимостью и безопасностью.

Современным открытым стандартом является «XLSX» (Office Open XML Spreadsheet, стандарт ISO/IEC 29500). Технически файл .xlsx представляет собой ZIP-архив, содержащий набор взаимосвязанных XML-файлов и папок. Например, xl/worksheets/sheet1.xml хранит данные ячеек, xl/sharedStrings.xml --- пул всех текстовых строк для оптимизации, xl/styles.xml --- описание стилей. Такая архитектура обеспечивает: открытость и хорошую совместимость между разными приложениями (LibreOffice, Google Sheets), структурную целостность, лучшую безопасность по сравнению с бинарными форматами и относительную компактность за счёт сжатия ZIP. Аналогичным открытым стандартом, используемым в бесплатных офисных пакетах (LibreOffice), является «ODS» (OpenDocument Spreadsheet, стандарт OASIS ISO/IEC 26300), также основанный на XML и ZIP \cite{ecma_office_open_xml_2016, quadratic_common_file_formats_2025}.

\ManualSubsubsection{Специализированные высокопроизводительные форматы для больших данных (Parquet, Avro)}

Обработка эксабайтов данных в распределённых системах (Hadoop, Spark) потребовала создания специализированных форматов, оптимизирующих скорость запросов и степень сжатия. «Apache Parquet» --- это «колоночный» бинарный формат. Вместо хранения данных построчно, он сохраняет значения каждого столбца вместе. Это позволяет системе при запросе «посчитать среднюю зарплату» читать с диска только столбец «Зарплата», игнорируя все остальные, что резко снижает ввод-вывод. Parquet применяет эффективные схемы кодирования (RLE, Dictionary) и сжатия (Snappy, GZIP) к однотипным данным в столбце, обеспечивая очень высокую степень компрессии. Также он поддерживает «проталкивание предикатов» --- фильтрацию целых блоков строк на основе метаданных (минимальное/максимальное значение в столбце), даже без их чтения \cite{apache_parquet_format}.

«Apache Avro», в отличие от Parquet, является «строчным» бинарным форматом. Его ключевая особенность --- хранение схемы данных в формате JSON прямо вместе с данными (обычно в начале файла). Это обеспечивает очень компактную сериализацию, высокую скорость чтения/записи и, что критически важно, безопасную «эволюцию схемы»: возможность добавлять новые поля или удалять старые при сохранении совместимости со старыми и новыми версиями приложений. Avro широко используется для сериализации сообщений в системах потоковой обработки данных, таких как Apache Kafka \cite{apache_parquet_format}.

\newpage
\refstepcounter{section}
\addcontentsline{toc}{section}{\protect\numberline{\thesection}Анализ и визуализация данных средствами табличных приложений}
\begin{center}\textbf{\thesection\ АНАЛИЗ И ВИЗУАЛИЗАЦИЯ ДАННЫХ СРЕДСТВАМИ ТАБЛИЧНЫХ ПРИЛОЖЕНИЙ}\end{center}
\vspace*{-1.5em}
\ManualSubsection{Инструментарий для многомерного анализа и исследования данных}

Современные табличные процессоры эволюционировали в полноценные аналитические платформы. Начальный этап анализа --- «упорядочивание и фильтрация». Многоуровневая сортировка позволяет организовать данные по нескольким критериям (например, сначала по региону по убыванию выручки, затем по менеджеру). Расширенные фильтры, включая фильтрацию по цвету, значению и сложным условиям («И»/«ИЛИ»), позволяют быстро выделять релевантные подмножества данных для детального изучения \cite{lakshmi_excel_analysis_2024}.

Ядром аналитического инструментария являются «функции». Помимо базовых математических и статистических (СУММ, СРЗНАЧ, СТАНДОТКЛОН), важную роль играют функции для анализа взаимосвязей (КОРРЕЛ --- корреляция, ЛИНЕЙН --- линейная регрессия) и прогнозирования (ТЕНДЕНЦИЯ, ПРЕДСКАЗ). Революцию в подходе к анализу произвели «динамические массивы», представленные в Excel 365. Функции типа ФИЛЬТР(), СОРТ(), УНИК() возвращают не одно значение, а сразу весь массив результатов, который автоматически «разливается» (spill) в соседние ячейки. Это позволяет создавать динамические отчёты, которые автоматически обновляются и меняют размер при изменении исходных данных без необходимости ручного копирования формул \cite{winston_excel_data_analysis_2016, probstein_data_analytics_excel365_2025}.

Наиболее мощным инструментом для сводного анализа по праву считаются «сводные таблицы» (Pivot Tables). Они позволяют, без написания формул, путём простого перетаскивания полей, мгновенно агрегировать (суммировать, усреднять, считать) большие массивы данных, группировать их по категориям, выполнять срезы по нескольким измерениям и рассчитывать доли. Интерактивные элементы управления --- «временные шкалы» для фильтрации по датам и «срезы» (Slicers) с кнопками для выбора значений --- превращают статический отчёт в интерактивный дашборд, с помощью которого можно исследовать данные «на лету» \cite{probstein_data_analytics_excel365_2025}.

\ManualSubsection{Принципы и процесс создания эффективных визуализаций}

Визуализация --- это искусство и наука преобразования чисел в инсайты. Первый и самый важный шаг --- «выбор типа диаграммы», который определяется вопросом, на который нужно ответить, и природой данных. Для сравнения величин по категориям используют столбчатые или линейчатые диаграммы; для демонстрации тренда во времени --- линейные; для отображения структуры целого (долей) --- круговые или кольцевые (хотя их применение часто критикуют за неточность восприятия); для исследования взаимосвязи между двумя числовыми переменными --- точечные; для отображения части от целого с накоплением --- каскадные; для воронки продаж --- воронкообразные. Специальные типы, такие как тепловые карты, эффективны для визуализации матриц данных и выявления «горячих» зон \cite{alexander_excel_charts_2016, intellezy_excel_vs_sheets_2024}.

Процесс создания диаграммы включает несколько этапов: 1. Подготовка и выбор корректного диапазона данных. 2. Вставка диаграммы нужного типа через интерфейс. 3. Детальная настройка: добавление и форматирование заголовков осей, легенды, меток данных, линий сетки. 4. «Оформление», подчиняющееся принципам информационного дизайна: «минимализм» (удаление всех нефункциональных элементов «чартджанка» --- 3D-эффектов, излишней штриховки), «создание визуальной иерархии» (самый важный элемент должен привлекать внимание первым), «целостность» (единый стиль всех элементов на дашборде) \cite{alexander_excel_charts_2016}.

Особое внимание уделяется «работе с цветом». Цвет должен использоваться семантически (красный --- опасность/убыток, зелёный --- безопасность/прибыль), а палитра должна быть ограниченной (5-7 основных цветов) и доступной для людей с дальтонизмом. Критически важно «не искажать данные»: масштаб осей должен начинаться с нуля для столбчатых диаграмм, нельзя сравнивать несопоставимые ряды на одном графике без вторичной оси. Интерактивность, добавляемая с помощью срезов сводных таблиц, позволяет пользователю самому выбирать параметры отображения, превращая статичный график в инструмент исследования \cite{alexander_excel_charts_2016}.

\newpage
\refstepcounter{section}
\addcontentsline{toc}{section}{\protect\numberline{\thesection}Автоматизация расчетов и использование формул и функций}
\begin{center}\textbf{\thesection\ АВТОМАТИЗАЦИЯ РАСЧЕТОВ И ИСПОЛЬЗОВАНИЕ ФОРМУЛ И ФУНКЦИЙ}\end{center}
\vspace*{-1.5em}

\ManualSubsection{Формулы как основа автоматизированных вычислений}

Автоматизация расчетов в табличных процессорах представляет собой методологию обработки данных, при которой повторяющиеся и сложные вычислительные задачи выполняются программными средствами по заранее определённому алгоритму, что минимизирует ручной труд и исключает ошибки, вызванные «человеческим фактором». Технической основой этой автоматизации являются формулы --- выражения, начинающиеся со знака равенства (=), которые производят вычисления над операндами с помощью операторов \cite{microsoft2024, belenkiy2020}.

Операндами могут быть: числа, текстовые константы, ссылки на ячейки (например, A1) или диапазоны (например, B2:B10). Использование ссылок --- ключевой элемент динамичности: при изменении значения в исходной ячейке результат формулы пересчитывается автоматически. Операторы определяют тип операции: арифметические, операторы сравнения и текстовый оператор конкатенации. Сложные формулы строятся по правилам приоритета операций и с использованием круглых скобок для группировки \cite{saffr2019, habr2023}.

\ManualSubsection{Классификация и применение встроенных функций для расширения возможностей}

Функции --- это именованные, предопределённые формулы, выполняющие специализированные вычисления по заданным аргументам. Они существенно расширяют аналитические возможности табличного процессора, позволяя решать сложные задачи без программирования. Функции можно классифицировать по их назначению \cite{exceljet2024, googlesupport2024}:

\begin{itemize}[label=-]
\item Математические и статистические: СУММ(), СРЗНАЧ(), МИН(), МАКС(), СЧЁТ() (подсчёт чисел), ДИСП() (дисперсия).
\item Логические: ЕСЛИ(условие; значение\_если\_истина; значение\_если\_ложь) --- базовый инструмент для ветвления вычислений; И() и ИЛИ() для проверки множественных условий.
\item Текстовые: СЦЕПИТЬ() (или оператор \&) для объединения строк; ЛЕВСИМВ(), ПРАВСИМВ(), ПСТР() для извлечения подстрок; СЖПРОБЕЛЫ() для удаления лишних пробелов.
\item Для работы с датами и временем: СЕГОДНЯ() (возвращает текущую дату), ДАТА(), ДЕНЬ(), МЕСЯЦ(), ГОД().
\item Функции поиска и ссылок: ВПР(искомое\_значение; таблица; номер\_столбца; [интервальный\_просмотр]) --- одна из ключевых функций для поиска данных в таблице по ключу.
\end{itemize}

Мощь функций раскрывается при их комбинировании (вложении). Например, конструкция =ЕСЛИОШИБКА(ВПР(A2; $F$2:$G$100; 2; ЛОЖЬ); "Не найден") выполняет поиск значения и корректно обрабатывает ситуацию ошибки, выводя понятное сообщение. Таким образом, освоение библиотеки функций и логики их сочетания является критически важным навыком, превращающим табличный процессор в мощный вычислительный инструмент для анализа данных любой сложности \cite{petrov2022, belenkiy2020}.


\newpage
\refstepcounter{section}
\addcontentsline{toc}{section}{\protect\numberline{\thesection}Основы подготовки табличных данных для последующего анализа}
\begin{center}\textbf{\thesection\ ОСНОВЫ ПОДГОТОВКИ ТАБЛИЧНЫХ ДАННЫХ ДЛЯ ПОСЛЕДУЮЩЕГО АНАЛИЗА}\end{center}
\vspace*{-1.5em}

\ManualSubsection{Принципы «чистых» данных и этапы подготовки}

Качество любого анализа напрямую зависит от качества исходных данных. Подготовка данных (data preprocessing) --- это комплексный процесс преобразования «сырых», неструктурированных данных в «чистый», пригодный для анализа формат. Его основная цель --- обеспечить точность, полноту, непротиворечивость и согласованность данных, что служит фундаментом для получения достоверных результатов и принятия обоснованных решений \cite{mckinney2017python, anderson2015creating}.

Процесс подготовки является итерационным и включает несколько ключевых этапов \cite{exceljet2024clean, habr2021datacleaning}:

\begin{enumerate}
    \item Сбор и первичный аудит: Консолидация данных из различных источников (файлы CSV, базы данных, веб-формы) и первоначальная оценка объёма, структуры и очевидных проблем (пропуски, аномалии);
    \item Очистка и стандартизация: Наиболее трудоёмкий этап, включающий:
        \begin{itemize}[label=-]
            \item Удаление дубликатов;
            \item Обработку пропущенных значений: анализ причин, решение об удалении строк или заполнении (константой, средним, медианой, модой);
            \item Исправление ошибок и опечаток: поиск орфографических ошибок, аномальных выбросов в числовых данных;
            \item Стандартизацию форматов: приведение дат к единому формату, текста --- к единому регистру, устранение лишних пробелов.
        \end{itemize}
    \item Преобразование и обогащение: Создание новых признаков на основе существующих (например, выделение дня недели из даты, расчёт возраста), разделение столбцов, агрегация данных;
    \item Верификация и сохранение: Финальная проверка корректности всех преобразований (например, с помощью сводных таблиц) и сохранение очищенного набора данных в новом файле или на новом листе.
\end{enumerate}

\ManualSubsection{Инструменты для очистки и трансформации данных в табличных процессорах}

Современные табличные приложения предлагают богатый набор инструментов для автоматизации рутинных задач очистки \cite{microsoft2024powerquery, googlesupport2024sheets}:

\begin{itemize}[label=-]
    \item «Удалить дубликаты»: Встроенный инструмент для быстрого поиска и удаления полностью совпадающих строк;
    \item «Текст по столбцам»: Мощный мастер для разделения данных, объединённых в одном столбце, по указанному разделителю (запятая, табуляция) или фиксированной ширине;
    \item «Найти и заменить» (Ctrl+H): Для массового исправления ошибок в написании, замены символов или стандартизации терминов;
    \item Проверка данных (Data Validation): Позволяет задать правила для вводимых значений (список, диапазон, дата), предотвращая появление некорректных данных на этапе их ввода;
    \item Специальные текстовые функции: СЖПРОБЕЛЫ(), ПРОПИСН(), СТРОЧН(), ПРОПНАЧ();
    \item Логические функции ЕСЛИ() и ЕСЛИОШИБКА(): Для условного преобразования данных и обработки ошибок в формулах.
\end{itemize}

Для сложных и повторяющихся задач преобразования данных в Excel существует надстройка Power Query (Get \& Transform). Это визуальная среда, позволяющая создавать воспроизводимые конвейеры (запросы) для загрузки данных из множества источников, их многоэтапной очистки, трансформации и загрузки в таблицу. Power Query запоминает все выполненные шаги, что позволяет одним кликом обновить весь процесс при появлении новых исходных данных, что критически важно для построения автоматизированных отчётов и дашбордов \cite{microsoft2024powerquery, dasu2003exploratory}.

\newpage
\refstepcounter{section}
\addcontentsline{toc}{section}{\protect\numberline{\thesection}Применение табличных инструментов в учебной и исследовательской работе}
\begin{center}\textbf{\thesection\ ПРИМЕНЕНИЕ ТАБЛИЧНЫХ ИНСТРУМЕНТОВ В УЧЕБНОЙ И ИССЛЕДОВАТЕЛЬСКОЙ РАБОТЕ}\end{center}
\vspace*{-1.5em}

\ManualSubsection{Роль в организации учебного процесса и научных исследований}

Табличные процессоры прочно вошли в академическую среду, став универсальным инструментом, значительно повышающим эффективность как учебной, так и научно-исследовательской деятельности. В образовательном процессе они используются для решения широкого спектра административных и учебных задач: от составления расписаний и планирования индивидуальных образовательных траекторий до контроля успеваемости через ведение электронных журналов с автоматическим расчётом среднего балла. Студенты используют таблицы для систематизации конспектов, управления библиографией для курсовых работ и планирования этапов своих проектов \cite{ivanova2023research, sidorov2022education}.

В научных исследованиях роль табличных инструментов ещё более значима. Они служат первичной средой для ввода, систематизации и предварительной обработки экспериментальных данных. Исследователь может непосредственно в таблицу заносить результаты измерений, наблюдений или анкетирования, после чего, не переходя в другое ПО, проводить их статистический анализ, строить графики и формировать прототипы таблиц для будущих публикаций. Это обеспечивает целостность и прозрачность исследовательского цикла на начальных этапах \cite{stanford2024research, walkenbach2019excel}.

\ManualSubsection{Специализированные возможности для анализа и визуализации научных результатов}

Для поддержки исследовательских задач табличные процессоры предлагают продвинутый аналитический функционал, выходящий за рамки базовой арифметики \cite{winston_excel_data_analysis_2016, albright2021analytics}:

\begin{itemize}[label=-]
    \item Статистические функции: Помимо описательной статистики (СРЗНАЧ, СТАНДОТКЛОН, ДИСП), доступны функции для проверки гипотез и анализа взаимосвязей, такие как КОРРЕЛ (коэффициент корреляции Пирсона) и ЛИНЕЙН (расчёт параметров линейной регрессии);
    \item Инструменты анализа «Что-Если»: «Подбор параметра» позволяет найти входное значение, необходимое для получения желаемого результата. «Поиск решения» решает задачи оптимизации с ограничениями (линейного и нелинейного программирования), что востребовано в экономических и инженерных расчётах;
    \item Пакет анализа (Analysis ToolPak): Надстройка в Excel, предоставляющая доступ к сложным статистическим методам: дисперсионному анализу (ANOVA), гистограммам, генерации случайных чисел, тестам на равенство средних (t-тест) и другим.
\end{itemize}

Визуализация результатов --- неотъемлемая часть научной коммуникации. Табличные процессоры позволяют создавать графики публикационного качества, соответствующие строгим академическим стандартам. Исследователь может построить точечную диаграмму с линией тренда для демонстрации корреляции, гистограмму для отображения распределения частот или комбинированный график для сравнения нескольких рядов данных. Условное форматирование помогает визуально выделять аномалии или тренды прямо в таблице с данными. Возможность создания интерактивных дашбордов на основе сводных таблиц и диаграмм со срезами позволяет как самому исследователю глубже изучать данные, так и эффектно представлять результаты коллегам или на конференциях \cite{alexander_excel_charts_2016, bullen2020data}.

Таким образом, владение продвинутыми возможностями табличных инструментов формирует ключевую цифровую компетенцию современного студента и исследователя, экономя время, повышая точность вычислений и качество визуального представления результатов, что в итоге способствует повышению эффективности и качества научно-исследовательской работы в целом \cite{ivanova2023research, sidorov2022education}.

\newpage
\section*{}
\vspace*{-2.5em}
\begin{center}\textbf{ЗАКЛЮЧЕНИЕ}\end{center}
\phantomsection
\addcontentsline{toc}{section}{Заключение}

Технологии создания и обработки табличных данных представляют собой целостный и логически связанный комплекс методов и инструментов, преобразующий необработанную информацию в стратегическое знание. Исследование показало, что эффективность работы с данными в конечном итоге определяется не только владением конкретными функциями программ, но и пониманием системных принципов, лежащих в основе каждого этапа этого процесса.

Фундаментом всей последующей работы являются принципы «чистых данных» и корректная первоначальная структуризация. Создание таблицы, где каждая строка — это наблюдение, а каждый столбец — переменная, формирует предсказуемую и машиночитаемую структуру, открывающую доступ к мощным возможностям сортировки, фильтрации и, что наиболее важно, сводного анализа. Пренебрежение этими принципами на старте делает любые попытки сложного анализа трудоёмкими и часто безуспешными.

Выбор технологии хранения и обработки на каждом этапе должен быть осознанным. Универсальные текстовые форматы (CSV, JSON) служат для надёжного обмена, открытые стандарты офисных пакетов (XLSX) — для комплексных документов, а специализированные колоночные форматы (Parquet) — для работы с большими данными. Понимание внутреннего устройства и сильных сторон каждого формата позволяет выстроить эффективный конвейер работы с информацией, минимизируя издержки на преобразования и потери данных.

Современные табличные процессоры предоставляют исчерпывающий инструментарий для превращения данных в инсайты. Динамические массивы и сводные таблицы позволяют гибко и наглядно исследовать данные, а правильно построенная визуализация, следующая принципам ясности и честности, делает полученные выводы убедительными и доступными для восприятия. При этом автоматизация рутинных расчётов с помощью формул и функций, а также применение инструментов для воспроизводимой очистки данных (Power Query) переводят работу из режима ручного труда в режим интеллектуального моделирования и анализа.

Таким образом, технологии работы с табличными данными — это не набор разрозненных приёмов, а взаимосвязанная экосистема. Комплексное владение этой экосистемой — от грамотного структурирования и выбора формата до продвинутого анализа и визуализации — формирует ключевую компетенцию специалиста в любой области, где решение основывается на данных. Это позволяет не только эффективно решать текущие задачи, но и создавать масштабируемые, воспроизводимые и надёжные процессы анализа, что является неотъемлемым условием успеха в условиях цифровой трансформации.

\newpage
\begin{center}\textbf{СПИСОК СОКРАЩЕНИЙ}\end{center}
\phantomsection
\addcontentsline{toc}{section}{Список сокращений}

\vspace{1em}

{\small
\noindent
\begin{tabular}{|p{0.15\textwidth}|p{0.80\textwidth}|}
\hline
\textbf{Сокращение} & \multicolumn{1}{c|}{\textbf{Расшифровка}} \\
\noalign{\hrule height 1.2pt}
\hline
CSV & Comma-Separated Values (значения, разделённые запятыми) \\
\hline
TSV & Tab-Separated Values (значения, разделённые табуляцией) \\
\hline
JSON & JavaScript Object Notation (обозначение объектов JavaScript) \\
\hline
XML & eXtensible Markup Language (расширяемый язык разметки) \\
\hline
XLS & Microsoft Excel Binary File Format (двоичный формат файлов Excel) \\
\hline
XLSX & Office Open XML Spreadsheet (электронная таблица Office Open XML) \\
\hline
ODS & OpenDocument Spreadsheet (электронная таблица OpenDocument) \\
\hline
Parquet & Apache Parquet (колоночный формат хранения данных) \\
\hline
Avro & Apache Avro (формат сериализации данных) \\
\hline
API & Application Programming Interface (интерфейс прикладного программирования) \\
\hline
ETL & Extract, Transform, Load (извлечение, преобразование, загрузка) \\
\hline
SQL & Structured Query Language (язык структурированных запросов) \\
\hline
UI & User Interface (пользовательский интерфейс) \\
\hline
UX & User Experience (пользовательский опыт) \\
\hline
BI & Business Intelligence (бизнес-аналитика) \\
\hline
VBA & Visual Basic for Applications (Visual Basic для приложений) \\
\hline
HTML & HyperText Markup Language (язык гипертекстовой разметки) \\
\hline
RDBMS & Relational Database Management System (реляционная система управления базами данных) \\
\hline
OLAP & Online Analytical Processing (оперативная аналитическая обработка) \\
\hline
OLTP & Online Transaction Processing (оперативная обработка транзакций) \\
\hline
SOA & Service-Oriented Architecture (сервис-ориентированная архитектура) \\
\hline
YAML & YAML Ain't Markup Language (YAML не является языком разметки) \\
\hline
\end{tabular}
}

\clearpage
\nocite{*}
\printbibliography[heading=bibliography]

\end{document}
